\chapter{Глобальный носитель и принуждение нетривиального сектора}

Предыдущая глава показала, что фаза полной обменной Real-пары не выбирает
класс пятнадцать относительно нулевого класса. Теперь проверяется более
сильная возможность: может ли сам глобальный носитель перехода исключить
тривиальный сектор ещё до построения частицы и её действия.

Проверка разделена на два утверждения. Первое является чистой теоремой о
главных окружностных расслоениях над $S^2$. Второе спрашивает, доказано ли
в проекте, что физическая сфера дефекта и физический фазовый носитель
действительно совпадают с этим расслоением.

\section{Классификация окружностных носителей}

Главные $U(1)$-расслоения
\begin{equation}
 U(1)\longrightarrow P_n\longrightarrow S^2
\end{equation}
классифицируются целым числом
\begin{equation}
 n=\langle c_1(P_n),[S^2]\rangle\in\mathbb Z.
 \label{eq:v5-carrier-chern-number}
\end{equation}
В длинной точной гомотопической последовательности граничное отображение
\begin{equation}
 \partial:\pi_2(S^2)\simeq\mathbb Z
 \longrightarrow
 \pi_1(U(1))\simeq\mathbb Z
\end{equation}
является умножением на $n$. Поэтому
\begin{equation}
 \pi_1(P_n)\simeq
 \begin{cases}
  \mathbb Z,&n=0,\\
  \mathbb Z/|n|\mathbb Z,&n\ne0.
 \end{cases}
 \label{eq:v5-carrier-fundamental-group}
\end{equation}

В стандартной классификации полное пространство имеет вид
\begin{equation}
 P_0\simeq S^2\times S^1,
 \qquad
 P_n\simeq L(|n|,1)\quad(n\ne0).
 \label{eq:v5-carrier-lens-spaces}
\end{equation}
В частности,
\begin{equation}
 P_{\pm1}\simeq S^3,
 \qquad
 P_{\pm2}\simeq L(2,1)\simeq\mathbb{RP}^3.
 \label{eq:v5-carrier-special-cases}
\end{equation}
Эти факты являются стандартными; в физической литературе та же семья
записывается как $M(0,p)=L(p,1)$, а случай $p=0$ --- как
$S^2\times S^1$. См. \path{arXiv:hep-th/0601068},
\path{arXiv:hep-th/0306062}, \path{arXiv:hep-th/0409073} и
\path{arXiv:0803.3241}.

\begin{theorem}[Топологическое принуждение генератора]
Если физический носитель переходной фазы задан как главное
$U(1)$-расслоение над $S^2$ с полным пространством $S^3$, то
\begin{equation}
 |c_1|=1.
\end{equation}
Тривиальный класс $c_1=0$ невозможен, поскольку его полное пространство
равно $S^2\times S^1$ и имеет фундаментальную группу $\mathbb Z$, тогда
как $S^3$ односвязна. Если вместо спинорного носителя фиксировано
$\mathbb{RP}^3$, то вынужден класс $|c_1|=2$.
\end{theorem}

Таким образом, после фиксации полного пространства вопрос о ненулевом
классе уже не является динамическим выбором вакуума. Остаётся только знак,
то есть ориентация линии.

\section{Следствие для коэффициентного класса пятнадцать}

Пусть $q_0$ --- уже построенный коэффициентный проектор ранга пятнадцать.
Внешнее произведение хопфова генератора с $[q_0]$ даёт
\begin{equation}
 \langle c_1(L)\,[q_0],[S^2]\rangle
 =15n.
 \label{eq:v5-carrier-coefficient-charge}
\end{equation}
Поэтому спинорный носитель $S^3$ принуждает ориентированные классы
$+15$ или $-15$, а нулевой носитель дал бы класс ноль. Для векторного
носителя $\mathbb{RP}^3$ соответствующая комплексная величина равна
$\pm30$; это не тот же представитель, что доказанный обменный класс
пятнадцать.

Следовательно, если физический носитель уже отождествлён с хопфовым
$S^3$, то ранее доказанный минимальный дефицит
\begin{equation}
 \frac{15}{105}=\frac17
\end{equation}
становится не только неизбежным внутри выбранного сектора: сам ненулевой
сектор также становится топологически обязательным.

\section{Что фиксирует соответствие Мориты}

Ориентационный функтор предыдущих глав имеет вид
\begin{equation}
 \mathcal T_L(X)=X\otimes L^{\deg X},
 \qquad
 \deg E=+1,
 \qquad
 \deg E^*=-1.
 \label{eq:v5-carrier-morita-functor}
\end{equation}
Он строго связывает прямую и обратную стрелки с $L$ и $L^*$, сохраняет
композицию и фиксирует относительный знак классов Черна.

Но формула \eqref{eq:v5-carrier-morita-functor} корректна для любой
комплексной линии $L$, включая тривиальную. Группа Пикара сферы равна
\begin{equation}
 \operatorname{Pic}(S^2)\simeq H^2(S^2,\mathbb Z)\simeq\mathbb Z,
\end{equation}
и нейтральный элемент $0$ является допустимой линией. Моритова
мультипликативность не выделяет генератор $1$ среди всех элементов этой
группы. Связь линий, групп Пикара и моритовых самоэквивалентностей
обсуждается, например, в \path{arXiv:math/0105001} и
\path{arXiv:math/0304048}.

\begin{nogocriterion}[Градуировка не выбирает класс Черна]
Из степеней $\deg E=+1$ и $\deg E^*=-1$ следует правило
$L\leftrightarrow L^*$ после выбора $L$, но не следует
$c_1(L)=1$. Подстановка тривиальной линии сохраняет все функториальные
тождеcтва. Поэтому моритовская архитектура сама по себе не исключает
нулевой сектор.
\end{nogocriterion}

\section{Сверка с уже доказанным багажом проекта}

В ранней геометрии действительно присутствуют оба глобальных объекта:
\begin{equation}
 S^3\simeq\operatorname{Spin}(3)\simeq SU(2),
 \qquad
 \mathbb{RP}^3\simeq SO(3).
\end{equation}
Кроме того, сфера ориентированных направлений имеет каноническое описание
\begin{equation}
 S^2\simeq\operatorname{Spin}(3)/\operatorname{Spin}(2)
 \simeq SU(2)/U(1),
 \label{eq:v5-spin-direction-sphere}
\end{equation}
и соответствующее расслоение $SU(2)\to S^2$ является хопфовым.

Однако проектная сфера точечного дефекта была получена иначе. Первичным
объектом там является неориентированная семейная ось
\begin{equation}
 P=\mathbf n\mathbf n^{\mathsf T}\in\mathbb{RP}^2.
\end{equation}
Для единичного класса потребовался условный подъём к ориентированной сфере
$S^2$ и массе $\mathbf n\cdot\boldsymbol\sigma$. Предыдущие строгие гейты
также установили, что пространственное $SO(3)$-действие не выводится из
одного следа $M_{35}$, а требуемого полуцелого семейного модуля внутри
$H_{15}/M_{35}$ нет.

Поэтому пока не доказана коммутативная идентификация
\begin{equation}
 \boxed{
 S^2_{\rm defect}
 \;=\;
 \operatorname{Spin}(3)/\operatorname{Spin}(2)
 }
 \label{eq:v5-missing-defect-spin-identification}
\end{equation}
вместе с утверждением, что физический фазовый носитель над этой сферой
есть именно спин-кадровое расслоение $S^3\to S^2$, а не независимо
назначенная линия.

\begin{theorem}[Условное исключение нулевого сектора]
Если будет доказано отождествление
\eqref{eq:v5-missing-defect-spin-identification} и физическая линия
перехода окажется ассоциированной линией канонического расслоения
\begin{equation}
 \operatorname{Spin}(2)\longrightarrow
 \operatorname{Spin}(3)\longrightarrow S^2_{\rm defect},
\end{equation}
то $|c_1|=1$, коэффициентный класс равен $\pm15$, а нулевой сектор
топологически невозможен. В текущей версии V само это отождествление не
выведено; поэтому исключение нулевого сектора остаётся условным.
\end{theorem}

\section{Статус и критерий остановки}

Получен не новый свободный механизм, а точный логический разрез:
\begin{equation}
 \boxed{
 \begin{gathered}
 S^3\to S^2\ \Longrightarrow\ |c_1|=1
 \quad\text{--- строго},\\
 |c_1|=1\ \Longrightarrow\ |15c_1|=15
 \quad\text{--- строго},\\
 S^2_{\rm defect}=\operatorname{Spin}(3)/\operatorname{Spin}(2)
 \quad\text{--- не доказано}.
 \end{gathered}}
\end{equation}

Нельзя продолжать, просто объявляя раннее $S^3$ физическим носителем
позднего дефекта. Следующий и единственный содержательный тест должен
строить отображение от доказанного проекторного дефекта к сфере
пространственных спин-направлений и проверять его совместимость с
моритовскими стрелками, Real-обменом и коэффициентным проектором $q_0$.
Если такого естественного отображения нет, глобальный путь выбора сектора
закрывается. Если оно существует, сектор пятнадцать становится
топологически вынужденным без фазовой или энергетической подгонки.

Следующий гейт:
\path{version5_spin_cover_defect_sphere_bridge_gate.tex}.

\begin{protocol}
\begin{enumerate}
 \item классификация $U(1)$-расслоений над $S^2$: \textbf{$c_1=n\in\mathbb Z$};
 \item тривиальный носитель: \textbf{$S^2\times S^1$};
 \item полный носитель $S^3$: \textbf{$|c_1|=1$};
 \item полный носитель $\mathbb{RP}^3$: \textbf{$|c_1|=2$};
 \item коэффициентный заряд для $S^3$: \textbf{$\pm15$};
 \item совместимость тривиальной линии с моритовым функтором:
 \textbf{да};
 \item выбор $|c_1|=1$ одной градуировкой Мориты: \textbf{нет};
 \item проекторная сфера дефекта: \textbf{$\mathbb{RP}^2$};
 \item ориентированный подъём к $S^2$: \textbf{условен};
 \item отождествление со сферой спин-направлений: \textbf{не доказано};
 \item нулевой сектор при доказанном spin-cover bridge: \textbf{исключён};
 \item нулевой сектор в текущей версии V: \textbf{ещё не исключён};
 \item следующий гейт: \textbf{мост спин-накрытия к сфере дефекта}.
\end{enumerate}
\end{protocol}

Машинный сертификат сохранён в
\path{s2t_v5_global_carrier_forced_nontrivial_sector_gate_results.json}.